\section{Exact differences in the small-tail band}\label{sec:band}

We compute the entire band
\[
\begin{gathered}
 26\le b\le149,\qquad
 \lambda=(6b-B-|\eta|,B,\eta),\\
 \eta_1\ge\cdots\ge\eta_4\ge0,\quad
 |\eta|\le27,\quad
 \eta_1\le B\le\left\lfloor\frac{6b-|\eta|}{2}\right\rfloor .
\end{gathered}
\]
In particular it contains every flagged partition of
Verification~\ref{ver:certificates}.
For each fixed tail \(\eta\), we will construct seven Laurent series
independent of \(b\). Their coefficients at
\(B,B-b,\ldots,B-6b\) will sum to \(720(Q(\lambda)-F(\lambda))\).
The same seven rows therefore serve every degree and every second part
in the displayed range.
Write \(\mathcal Q_b=h_b[h_6]\), \(\mathcal F_b=h_6[h_b]\), and set
\[
 \mathcal C(q,y)=(1-q)\prod_{\ell=1}^4(1-y_\ell)(1-y_\ell/q),\quad
 \mathcal A(y)=\prod_{\ell<h}(1-y_h/y_\ell),\quad
 D(q)=\prod_{h=1}^6(1-q^h).
\]
We first expand the characters and geometric products in the tail
variables \(y\), and then expand their rational coefficients in \(q\)
as Laurent series at \(q=0\). In particular, \(D(q)^{-1}\) has only
nonnegative powers.

\subsection{Coefficient extraction and truncation}
The normalized alternating Schur formula gives
\begin{equation}\label{eq:bandextract}
 Q(\lambda)-F(\lambda)
 =[q^By^\eta]\mathcal C(q,y)\mathcal A(y)
                       (\mathcal Q_b-\mathcal F_b)(1,q,y).
\end{equation}
Indeed in six variables it extracts the coefficient of \(x^\lambda\)
after multiplication by \(\prod_{i<j}(1-x_j/x_i)\). Distinct decreasing
shifted partitions cannot be permutations of each other, so only the
Schur term at \(\lambda\) contributes. All Laurent monomials have total
degree \(6b\); setting the first variable to one loses no information
about its exponent once the remaining five are fixed.

Give each \(y_\ell\) degree one and \(q\) degree zero.
The monomials of \(\mathcal A\) have total \(y\)-degree zero; those of
\(\mathcal C\) have nonnegative total \(y\)-degree. Character terms above
degree \(27\) therefore cannot contribute to \eqref{eq:bandextract}.
This uses total degree even when individual exponents in the alternant
are negative. Write \(\equiv_{27}\) for equality of coefficients of
nonnegative \(y\)-monomials through total degree \(27\), with coefficients
in \(\mathbb Q(q)\).

\subsection{Seven numerators independent of the degree}
For \(0\le j\le6\) define
\[
 L_j(q,y)=
 \prod_{\substack{\tau\in\NN^4,\ 1\le|\tau|\le6\\0\le u\le6-|\tau|}}
                         (1-y^\tau q^{u-j})^{-1}.
\]
Then
\begin{equation}\label{eq:targetband}
 \mathcal Q_b(1,q,y)\equiv_{27}
 \frac1{D(q)}\sum_{j=0}^6q^{jb}(-1)^j
                  e_j(q,q^2,\ldots,q^6)L_j(q,y).
\end{equation}
To prove this, fix a multiset of \(k\) inner monomials with nonzero tail.
The remaining \(b-k\) entries come from \(1,q,\ldots,q^6\), with
generating polynomial
\[
 h_{b-k}(1,q,\ldots,q^6)=
                    \prod_{h=1}^6\frac{1-q^{b-k+h}}{1-q^h}
                                      \quad(b-k\ge0).
\]
This Gaussian polynomial counts partitions in a \(6\)-by-\((b-k)\)
rectangle. In general their generating polynomial \(P_{r,s}\), with at
most \(r\) rows of length at most \(s\), satisfies
\(P_{r,s}=P_{r,s-1}+q^sP_{r-1,s}\), by removing a row of length \(s\)
when one is present, and \(P_{0,s}=P_{r,0}=1\).
The displayed product with \(r=b-k,s=6\) satisfies the same recurrence
and boundary values, proving the formula.
Expanding the numerator gives the \(e_j\) in \eqref{eq:targetband};
the factor \(q^{-jk}\) contributes \(q^{-j}\) per selected tail monomial,
and hence produces \(L_j\).
Within the cutoff \(k\le27\) and \(b\ge26\), so the only possible
negative value of \(b-k\) is \(-1\). Its numerator has the factor
\(1-q^0\) and vanishes, exactly as required. Thus the formula includes
the boundary \(b=26,k=27\); its cancellation is across the seven
\(j\)-terms.

Similarly, the binary complement of one tail monomial gives
\begin{equation}\label{eq:binarysource}
 h_b(1,q,y)\equiv_{27}
 \frac{\prod_\ell(1-y_\ell)^{-1}
       -q^{b+1}\prod_\ell(1-y_\ell/q)^{-1}}{1-q}.
\end{equation}
At tail degree \(t\), its coefficient is
\((1-q^{b+1-t})/(1-q)\). It counts the binary monomials when \(t\le b\)
and is zero at \(t=b+1\), the only extra possible degree.

For a cycle type \(I=(i_1,\ldots,i_v)\vdash6\), let
\[
 n_I=\frac{720}{\prod_{r\ge1}r^{m_r}m_r!},\qquad
 E_I(q)=\frac{D(q)}{\prod_{h=1}^v(1-q^{i_h})},
\]
where \(m_r\) counts cycles of length \(r\).
The quotient \(E_I\) is an integer polynomial of degree \(15\):
the multiplicity of the cyclotomic factor \(\Phi_d\) in its denominator
is the number of cycle lengths divisible by \(d\), which is at most
\(\lfloor6/d\rfloor\), its multiplicity in \(D\).
Apply \eqref{eq:binarysource} with variables raised to each cycle length
in \eqref{eq:cycle}. Together with \eqref{eq:targetband} this yields
\begin{equation}\label{eq:seven}
 720D(q)(\mathcal Q_b-\mathcal F_b)(1,q,y)
                         \equiv_{27}\sum_{j=0}^6q^{jb}G_j(q,y),
\end{equation}
where the seven numerators are explicitly
\begin{align}
G_j(q,y)={}&720(-1)^je_j(q,\ldots,q^6)L_j(q,y)\notag\\
 &-\sum_{I\vdash6}n_I
   \sum_{\substack{J\subseteq\{1,\ldots,v\}\\\sum_{h\in J}i_h=j}}
   (-1)^{|J|}q^j E_I(q)
   \prod_{\ell=1}^4\prod_{h=1}^v
       (1-y_\ell^{i_h}q^{-i_h\one_{h\in J}})^{-1}.
                                                     \label{eq:G}
\end{align}
Here \(\one_{h\in J}\) is the indicator of membership.
Every selected second term of \eqref{eq:binarysource} contributes
\(-q^{i_h(b+1)}\), explaining the sign and the residual \(q^j\)
after \(q^{jb}\) is factored out. Equal-length cycles still have distinct
positions in the subset sum. These formulas involve eleven cycle types
and seven values of \(j\), regardless of \(b\).

\subsection{Computing the numerators}
Write \(L_{j,a}=[y^a]L_j\). Euler differentiation of the product gives
\begin{equation}\label{eq:euler}
 L_{j,0}=1,\qquad
 |a|L_{j,a}=
 \sum_{\tau,u}\ \sum_{\substack{i\ge1\\i\tau\le a}}
              |\tau|q^{i(u-j)}L_{j,a-i\tau}\quad(|a|>0).
\end{equation}
For example, this follows by differentiating the logarithmic series
\(\sum_{i\ge1}(y^\tau q^{u-j})^i/i\). Its denominator \(i\) cancels
the extra \(i\) from differentiation. Every lookup has smaller total
tail degree, and the division by \(|a|\) is exact.

For the source product in \eqref{eq:G}, put
\[
 A_J(m)=[z^m]\prod_{h\in J}(1-z^{i_h})^{-1}.
\]
At a single tail-coordinate exponent \(m\), its Laurent coefficient is
\[
          \sum_{p=0}^m A_J(p)A_{\bar J}(m-p)q^{-p}.
\]
Multiply these four one-coordinate polynomials and then
\((-1)^{|J|}n_Iq^jE_I(q)\). These are the \texttt{ways} arrays and
convolutions of \texttt{source} in Appendix~\ref{app:band}.
The routines \texttt{target} and \texttt{gen\_source} implement
\eqref{eq:euler} and the cycle-type enumeration.

All raw arrays are symmetric in the four tail variables, so they need
only be stored at decreasing nonnegative exponent vectors; arbitrary
orders are looked up by sorting, without a sign.
To multiply by the cross factors in \(\mathcal C\), the coefficient
at a shift \(v\in\{0,1,2\}^4\) is
\[
                  (-1)^{n_1}q^{-n_2}(1+q^{-1})^{n_1},
       \qquad n_s=\#\{\ell:v_\ell=s\}.
\]
An in-place sweep in decreasing tail degree reads unmodified lower-degree
rows. A subsequent decreasing-exponent sweep multiplies by \(1-q\).
This is the routine \texttt{cross}.

The cross factors preserve symmetry in the four tail variables. To
extract a Schur coefficient, we now apply the alternating factor
\(\mathcal A\); the permutation signs enter at this step, rather than in the
sorted lookups.
Put \(\rho=(3,2,1,0)\), and let \(M_{j,a}=[y^a]\mathcal C G_j\),
with value zero at a vector with a negative coordinate.
The exact row produced by \texttt{schur\_row} is
\begin{equation}\label{eq:row}
 R_{j,\eta}(q)=\frac1{D(q)}
       \sum_{\pi\in\SSS_4}\sgn(\pi)
                     M_{j,\eta+\rho-\pi\rho}(q).
\end{equation}
This follows from
\(\mathcal A(y)=y^{-\rho}\sum_\pi\sgn(\pi)y^{\pi\rho}\).
Every accepted index in \eqref{eq:row} has nonnegative entries and total
degree \(|\eta|\le27\); in particular every coordinate is at most \(27\).
There are \(1908\) decreasing four-coordinate tails in this range.
Equations~\eqref{eq:bandextract} and \eqref{eq:seven} show that the final
integer tested is precisely
\begin{equation}\label{eq:finalinteger}
 \sum_{j=0}^6[q^{B-jb}]R_{j,\eta}(q)
       =720\bigl(Q(6b-B-|\eta|,B,\eta)-F(6b-B-|\eta|,B,\eta)\bigr).
\end{equation}
The notation on the right specifies partitions rather than gap vectors.

\subsection{The Laurent window}
At total tail degree \(t\), the coefficients and individual Euler
summands of \(L_j\) have \(q\)-support in \([-6t,5t]\).
The elementary factor adds at most \(21\). A cross term consuming
\(s\) tail units shifts down by at most \(s\), while its raw input has
degree \(t-s\); the factor \(1-q\) adds at most one.
Thus all target numerator summands lie in
\[
                      [-6t,5t+22].
\]
The source has raw support in \([-t,21]\) and final support in
\([-t,22]\), using \(\deg E_I=15\) and \(0\le j\le6\).
Alternation does not change total tail degree or \(q\)-support.
At \(t\le27\) all these supports are therefore contained in
\[
                       [-162,157]\subset[-168,161],
\]
the interval stored by the program. Its clipping operation discards
only zero terms. In particular all terms of the boundary cancellation
at \(b=26,k=27\) are retained separately.

Since \(D(0)=1\), its inverse is a power series with no negative
exponents. Starting from exponent \(-168\), six forward passes
\(a_n\gets a_n+a_{n-h}\), for \(h=1,\ldots,6\), divide by \(D\)
exactly. The numerator is zero above \(161\).
Expansion through \(447\) suffices because
\[
                         B-jb\le B\le3b\le447.
\]
Entries below \(-168\) are zero; every accessed nonnegative array index
is at most \(168+447=615\), within the \(616\) allocated entries.
Neither division nor the shift \(q^{jb}\) can bring a discarded high
exponent down to a requested one. The program converts to
arbitrary-precision integers before the denominator passes and uses them
for all final sums and divisibility tests.

Finally, \texttt{run\_thin} visits each of the \(1908\) tails once,
constructs its seven rows, and loops over the displayed ranges of
\(b\) and \(B\). For a fixed \(b,\eta\), the number of permitted
second parts is
\[
       3b-\left\lceil\frac{|\eta|}{2}\right\rceil-\eta_1+1.
\]
This is positive throughout the band. Summing it over the tails and
\(26\le b\le149\) gives \(57065668\) distinct partitions; there is
no dependence on which partitions the first program flagged.

\begin{verification}\label{ver:band}
The complete program in Appendix~\ref{app:band} checks \(57065668\)
coefficients in this band. Of these, \(177134\) are zero and none are
negative; every integer in \eqref{eq:finalinteger} is divisible by \(720\).
The grouped counts appear in Table~\ref{tab:band}.
\end{verification}
